\section{The constructive proof is already within reach}\label{ch09:sec:constructive}
Chapter~\ref{ch:matching} already contains almost everything we need for the sufficient direction. Recall the complete search of Section~\ref{ch08:sec:search}, started at an unmatched student $x$ and relative to a matching $M$. From a reached student, the search follows every edge that is not in $M$ to the task at its other end. When it reaches a task that is held by some student, it also reaches that task's owner. Each vertex is reached at most once, and the search remembers its predecessor, the vertex from which it was reached. There are finitely many vertices, so the search ends. Chapter~\ref{ch:matching} proved two facts about how it can end.
\begin{itemize}
\item If the search reaches an unused task, then following the predecessors back to $x$ gives an augmenting path. Flipping $M$ along this path gives a matching with one more edge (Lemma~\ref{thm:augment}).
\item If the search fails, that is, if it ends without reaching an unused task, let $A$ be the set of students it reached, $x$ included, and let $B$ be the set of tasks it reached. Then $N(A)=B$ and $|B|=|A|-1$ (Proposition~\ref{ch08:prop:shortage}). Hence $|N(A)|=|A|-1<|A|$.
\end{itemize}
In the second case, $A$ is a shortage set. Hall's condition says that no group of students is short of tasks, so under \eqref{eq:hall} the second case can never happen. That is the whole idea of the proof.

\begin{proof}[Constructive proof that Hall's condition is sufficient]
Assume that the graph satisfies \eqref{eq:hall}. Start with the empty matching. As long as some student is unmatched, carry out the following round. Choose an unmatched student $x$ and run the complete search from $x$. If the search failed, the set $A$ of students it reached would satisfy $|N(A)|<|A|$, contradicting \eqref{eq:hall}. So the search reaches an unused task, and flipping the matching along the resulting augmenting path increases its number of edges by one.

Different edges of a matching contain different students, so a matching has at most $|L|$ edges. The matching starts with $0$ edges and gains one edge in every round, so after $k$ rounds it has $k$ edges, and therefore $k\le|L|$. Thus there are at most $|L|$ rounds. The rounds stop only when no student is unmatched, so the final matching saturates $L$.
\end{proof}

If $L=\varnothing$, no round takes place at all. The empty matching already saturates $L$, because there is no student to cover.

Section~\ref{sec:matching-trace} shows both outcomes on a small example. With $N(a)=\{1,2\}$, $N(b)=\{1\}$ and $N(c)=\{2,3\}$, the graph satisfies Hall's condition (you can check all eight inequalities). The first two rounds there give $a$ task~$1$ and $c$ task~$2$. This is an unlucky start, because task~$1$ is the only task that $b$ accepts. The third round repairs it along the augmenting path $b$--$1$--$a$--$2$--$c$--$3$, exactly as the proof promises. When task~$3$ is removed from $c$'s options, Hall's condition fails for the whole class, which then accepts only tasks $1$ and~$2$. The search from $b$ fails, and it reports precisely the group $\{a,b,c\}$ where the condition fails.

Chapter~\ref{ch:matching} observed that the procedure, run on any finite bipartite graph, ends with a certificate. Either it produces a matching that saturates $L$, which can be checked pair by pair, or it produces a shortage set, which can be checked by counting. Hall's theorem adds that, for each graph, exactly one of the two kinds of certificate exists. If some matching saturates $L$, the necessary direction shows that there is no shortage set. If no matching saturates $L$, the sufficient direction shows that \eqref{eq:hall} fails, and any group where it fails is a shortage set. So the procedure can never mistake an unlucky sequence of early choices for a genuine impossibility. Unlucky choices are repaired along augmenting paths, and only a real shortage can stop the procedure.
